\chapter{The interfaces of a reproducible account}\label{ch:implementation}

A theorem names a mathematical object. An EBU system must preserve that object as observations become proposals, proposals become accepted actions and actions become records. This chapter specifies the theoretical interfaces that keep those meanings intact. It explains why a reproducible account can remain local, composable and open to correction without becoming a single opaque score.

\section{Every quantity has a role}
A physical stock, a reference, a rate, a finite quantity and an institutional balance can all be stored as numbers. Their meanings are nevertheless different. A stock is available material; a reference is a comparison condition; a rate acts over time; a finite quantity is the resulting amount; a balance belongs to a policy. Substituting one for another can produce plausible arithmetic while changing the question.

The account therefore preserves the type, unit and identity of each object. Stocks and their units travel with the cell map. References, scales and factor definitions travel with the potential version. Action directions and accepted quantities travel with the complete transformation. Time horizons, pending state and controller parameters travel with the response declaration.

This is more than careful naming. It is the mathematical condition for combining records. Two sites can add carrier flows only in compatible units and boundaries. Two consecutive EBU receipts telescope only when their states join and their field is held fixed, or when changes are separately recorded. A capacity rule can use a receipt without pretending that EBU is another stock of water.

\section{Permission, response and valuation}
Permission identifies the admissible transformations. The response map supplies the complete endpoint of an accepted transformation. Valuation compares that endpoint with the declared baseline. An action selector chooses among alternatives using a separately specified policy.

For an immediate lossless transfer, the response may simply be $x'=x+S_eq$. For a maintained service it may be the endpoint of a differential equation at $T_f$. For a shared source it may include a resolver that reduces competing requests. These are different response interfaces to the same finite definition.

The interfaces can be assembled operationally, but their responsibilities remain distinct. An evaluator that silently changes a requested quantity has also become a resolver. A permission rule that chooses the largest EBU has also become a selector. A recorder that changes the baseline to make a total close has changed the event it claims to describe.

A clear separation supports efficient improvement. A better route model can replace the old response while preserving the valuation definition. A fuller potential can be introduced with an explicit field version. A different access rule can use the same physical receipts. Each revision then has an identifiable purpose and effect.

\section{The event that the receipt describes}
A complete event record connects its identity, predecessor, accepted transformation, successor and field. It retains the service condition and the physical evidence relevant to those inputs. A joint event retains its child actions and the convention used to attribute the total.

For a dynamic service, the record also identifies the horizon and background comparison. The endpoint includes whatever remains pending at that time. The raw interval change and the action-relative change can both be recorded, but they are not interchangeable. Their difference is the no-action background term.

This lets a participant reconstruct the number from the same operands. It also makes disagreement specific. A mistaken meter reading changes an observed endpoint; a changed reference changes the potential; a disputed owner changes attribution. The original record remains visible when a correction is appended.

\section{Local evaluation follows exact support}
If $V=\sum_\alpha\phi_\alpha$, factors with unchanged inputs cancel from the finite difference. The evaluator therefore needs only the affected factors and their complete arguments. The efficiency follows from the theorem, not from an assumption that distant places never matter.

Response can enlarge the affected endpoint support. A dispatched action can reach a downstream compartment before the horizon, or alter a shared reserve through feedback. Its local calculation must include the factors changed by that complete response. Pending state cannot be dropped merely because it is outside the receiver's visible tank.

Where a composed factor depends on at most $r$ action labels, its higher differences vanish above order $r$. Where a common linear response feeds a quadratic field, the ceiling is two. These results let an implementation retain an exact smaller representation when its declarations satisfy the premises.

The practical benefit is less redundant work. A small action need not trigger an entire planetary recalculation. At the same time, the relevant shared resource is included when it matters. This combination is central to an account that can serve many small participants rather than only organisations able to maintain a large central model.

\section{Three equivalent calculations, one value}
The four-unit transfer from $(18,2)$ to $(14,6)$ has EBU twelve under the running field. Endpoint subtraction gives $16-4$. Quadratic expansion gives $4(4)-(1/4)4^2$. Integrating the changing field gives $\int_0^4(4-q/2)\,dq$. Their equality follows from the earlier laws.

These representations explain different parts of the same event. Endpoints establish the total; curvature explains why the initial slope overstates it; the path integral explains how the marginal changes. They are not three contributions to add.

The same discipline applies to interaction. A group total reconstructed from its Möbius coefficients equals the original table entry. The coefficients explain the total and must not be issued again as additional value. A common-path allocation or dividend allocation then divides the known total under its declared rule.

\section{History has an exact reconciliation}
For consecutive complete states under one field, the receipt total is $V(x_0)-V(x_m)$. External movement has its own line; a field revision has its own line; an attribution has its own sum. Each identity checks a particular connection between records.

This makes missing or duplicate claims visible. A positive restoration cannot be detached from a preceding negative action when the purpose is to report the complete net history. A split transfer cannot reuse the original baseline for every piece. An autonomous controller cannot create a new registered actor event with every oscillation.

The benefit is a history that can support both accountability and maintenance. An institution can recognise useful work once, preserve necessary burdens honestly and see where recurring resource use suggests a physical improvement.

\section{Approximation remains an explicit object}
A local Taylor expression may approximate a finite result; an estimated endpoint may approximate a physical state; a reduced coalition model may omit higher terms. Each approximation has a target and a remainder or error description.

For a truncated multilinear account, the remainder at a coalition is exactly the sum of omitted coefficients contained in that coalition. For smooth small actions, the derivative assumptions give the Taylor order. For uncertain table entries, covariance or deterministic bounds propagate through the signed contrast. The word approximation therefore need not mean an unexplained departure from the theory.

This gives a system a useful flexibility. It can use economical calculations where a controlled error is acceptable and exact finite formulas where the event is consequential or the simplification is available. The reason for the choice remains visible to the people relying on the result.

\section{A reproducible account remains an understandable account}
A reader should be able to move from the service being supplied to the physical change, from that change to its EBU, and from the EBU to any institutional consequence. These are the same connections developed throughout the book. Their implementation should preserve the explanation rather than replace it with an unexplained authority.

The next chapter assembles the declarations into one coherent theoretical model. Detailed software verification, measurement designs, simulation protocols and results belong to the later books; the present task is to specify exactly what those implementations and studies will be about.
